\usetheme{ugm1}
\usepackage{array}
\usepackage{booktabs}
\usepackage{tabularx}
\usepackage{amsmath}
\usepackage{fix-cm}
\usepackage{tikz}
\usetikzlibrary{positioning,fit,calc,arrows.meta}

\hypersetup{
    pdftitle={Perbandingan Terkontrol Prosedur Imperatif dan GitOps dalam Menghadapi Penyimpangan Konfigurasi pada Klaster Kubernetes},
    pdfauthor={Muhammad Argya Vityasy},
    pdfsubject={Sidang Skripsi},
}
\pdfstringdefDisableCommands{\def\translate#1{#1}}

\definecolor{gitopsblue}{RGB}{51,108,229}
\definecolor{imperativegray}{RGB}{102,102,102}
\definecolor{driftred}{RGB}{229,115,115}
\definecolor{okgreen}{RGB}{76,175,80}

% Judul lengkap harus muat di halaman sampul tema.
\setbeamerfont{title}{size=\fontsize{13pt}{14.2}\selectfont,family=\gamasans,series=\bfseries}

% -------------------------------------------------------
% Presentation metadata
% -------------------------------------------------------
\title{\texorpdfstring{\raggedright\hyphenpenalty=10000\exhyphenpenalty=10000 Perbandingan Terkontrol Prosedur Imperatif dan GitOps dalam Menghadapi Penyimpangan Konfigurasi pada Klaster Kubernetes}{Perbandingan Terkontrol Prosedur Imperatif dan GitOps dalam Menghadapi Penyimpangan Konfigurasi pada Klaster Kubernetes}}
\subtitle{Sidang Skripsi\texorpdfstring{\\}{} Studi Kasus InvenioRDM}
\author{Muhammad Argya Vityasy\texorpdfstring{\\[3pt]}{, }%
NIM: 23/522547/PA/22475\texorpdfstring{\\[2pt]}{, }%
Dosen Pembimbing: Guntur Budi Herwanto, S.Kom., M.Cs.\texorpdfstring{\\[2pt]}{, }%
Dosen Penguji 1: Dr.\ techn.\ Kabul Kurniawan, S.Kom., M.Cs.\texorpdfstring{\\[2pt]}{, }%
Dosen Penguji 2: Aina Musdholifah, S.Kom., M.Kom., Ph.D.}
\date{2026}

\graphicspath{{diagrams/}{../proposal/diagrams/}}
\newcommand{\fullslidediagram}[2][0.68\textheight]{%
    \includegraphics[width=0.96\textwidth,height=#1,keepaspectratio]{#2}%
}
\newcommand{\slidelead}[1]{%
    \begin{beamercolorbox}[wd=\textwidth, sep=0.20cm, center, rounded=true]{structure}
        \small\bfseries #1
    \end{beamercolorbox}
}
\newcommand{\smallnote}[1]{%
    \begin{beamercolorbox}[wd=\textwidth, sep=0.16cm, center, rounded=true]{frametitle}
        \scriptsize #1
    \end{beamercolorbox}
}
\setbeamertemplate{frametitle}{%
    \vspace*{0.52cm}\hspace*{-0.25cm}%
    \begin{beamercolorbox}[wd=0.82\textwidth, dp=0.12cm]{frametitle}
        \usebeamerfont{frametitle}\insertframetitle\strut
    \end{beamercolorbox}%
    \vspace*{0.10cm}%
}

\begin{document}

% =============================================
% 1. Title
% =============================================
\begin{frame}
    \titlepage
\end{frame}

% =============================================
% 2. Problem
% =============================================
{
\usebackgroundtemplate{\backgroundcontent}
\begin{frame}
    \frametitle{Configuration Drift Bukan Cacat Statis}
    \begin{columns}[T]
        \begin{column}{0.50\textwidth}
            \scriptsize
            \textbf{Drift} adalah ketidakcocokan \textit{actual state} klaster dengan deklarasi acuan,
            bukan penilaian atas kebenaran acuan.
            \vspace{0.06cm}

            \textbf{Dampaknya nyata:}
            \begin{itemize}\setlength{\itemsep}{1pt}
                \item 719 \textit{defect} konfigurasi pada 2.260 skrip Kubernetes;
                      533 di antaranya (74,1\%) menyebabkan \textit{crash}, operasi salah, atau \textit{outage} (Zhang dkk., 2026).
                \item \textit{Security misconfiguration} pada manifest \textit{open-source} berulang dalam 11 kategori (Rahman dkk., 2023).
            \end{itemize}
            \vspace{0.04cm}
            \textbf{Pertanyaannya:} bagaimana dua prosedur manajemen merespons drift yang sama?
        \end{column}
        \begin{column}{0.46\textwidth}
            \centering
            \includegraphics[width=\columnwidth,height=0.50\textheight,keepaspectratio]{motivation-drift}
        \end{column}
    \end{columns}
    \vspace{0.04cm}
    \slidelead{Drift mengubah kondisi operasional yang harus dipulihkan.}
\end{frame}
}

% =============================================
% 3. Gap
% =============================================
{
\usebackgroundtemplate{\backgroundcontent}
\begin{frame}
    \frametitle{Celah: Perbandingan Terkontrol pada Workload Nyata}
    \begin{columns}[T]
        \begin{column}{0.48\textwidth}
            \scriptsize
            \textbf{Perbandingan sudah ada}
            \begin{itemize}\setlength{\itemsep}{3pt}
                \item Shrestha dan Ali (2024) membandingkan GitOps dan Ansible pada waktu induksi, deteksi, serta pemulihan.
                \item Kagganti Nataraja (2025) menguji perbandingan terkontrol pada Minikube satu node.
                \item Matubber (2025) mengevaluasi GitOps dengan pembanding semi-otomatis.
            \end{itemize}
        \end{column}
        \begin{column}{0.48\textwidth}
            \scriptsize
            \textbf{Yang belum dipisahkan}
            \begin{itemize}\setlength{\itemsep}{3pt}
                \item Kesetaraan konfigurasi dari sinyal kesehatan yang dipakai sebagai \textit{guardrail}.
                \item Sinyal deteksi dari ketersediaan bukti dan pemulihan.
                \item Prosedur imperatif otomatis yang tetap, terpisah dari kecepatan operator.
            \end{itemize}
        \end{column}
    \end{columns}
    \vspace{0.10cm}
    \slidelead{Kebutuhan: \textit{workload} multi-service nyata, prosedur yang dapat diulang, dan metrik yang ditetapkan sebelum pengumpulan.}
    \vspace{0.06cm}
    \smallnote{Kontribusi bersifat terbatas: perbandingan pada \textit{apparatus} yang diuji, bukan klaim superioritas universal.}
\end{frame}
}

% =============================================
% 4. Research questions
% =============================================
{
\usebackgroundtemplate{\backgroundcontent}
\begin{frame}
    \frametitle{Rumusan Masalah: Satu Pertanyaan Waktu}
    \vspace{0.04cm}
    \begin{beamercolorbox}[wd=\textwidth, sep=0.18cm, rounded=true]{structure}
        \scriptsize
        Seberapa besar perbedaan waktu deteksi, waktu tersedianya bukti yang dapat diverifikasi,
        dan waktu pemulihan konfigurasi penuh antara prosedur imperatif otomatis dan GitOps
        pada setiap skenario drift, dan pada kondisi apa keunggulan waktu menyusut atau bertahan?
    \end{beamercolorbox}
    \vspace{0.10cm}
    \begin{columns}[T]
        \begin{column}{0.31\textwidth}
            \begin{beamercolorbox}[wd=\textwidth, sep=0.10cm, center, rounded=true]{frametitle}
                \scriptsize\bfseries (a) Deteksi
            \end{beamercolorbox}
            \vspace{0.04cm}
            \centering\scriptsize $t_d - t_0$
        \end{column}
        \begin{column}{0.31\textwidth}
            \begin{beamercolorbox}[wd=\textwidth, sep=0.10cm, center, rounded=true]{frametitle}
                \scriptsize\bfseries (b) Bukti
            \end{beamercolorbox}
            \vspace{0.04cm}
            \centering\scriptsize $t_e - t_0$
        \end{column}
        \begin{column}{0.31\textwidth}
            \begin{beamercolorbox}[wd=\textwidth, sep=0.10cm, center, rounded=true]{frametitle}
                \scriptsize\bfseries (c) Pemulihan
            \end{beamercolorbox}
            \vspace{0.04cm}
            \centering\scriptsize $t_r - t_0$
        \end{column}
    \end{columns}
    \vspace{0.12cm}
    \smallnote{``Pada kondisi apa'' dijawab oleh perbandingan per skenario, bukan subpertanyaan keempat. R1 = validitas pemulihan; kesehatan layanan = \textit{guardrail}. A, F, dan G dilaporkan terpisah.}
\end{frame}
}

% =============================================
% 5. Design
% =============================================
{
\usebackgroundtemplate{\backgroundcontent}
\begin{frame}
    \frametitle{Desain: Satu \textit{Stack}, Dua Mode Berurutan}
    \centering
    \includegraphics[width=0.98\textwidth,height=0.40\textheight,keepaspectratio]{design}
    \vspace{0.04cm}
    \begin{columns}[T]
        \begin{column}{0.48\textwidth}
            \scriptsize
            \textbf{Yang dikendalikan:} manifest, \textit{image}, instrumen, HPA, tujuh skenario, dan kondisi dasar.
        \end{column}
        \begin{column}{0.48\textwidth}
            \scriptsize
            \textbf{Yang membedakan:} mekanisme manajemen. ArgoCD aktif pada mode GitOps;
            ArgoCD ditangguhkan sementara dan pemeriksa \texttt{kubectl} bekerja pada mode imperatif.
        \end{column}
    \end{columns}
    \vspace{0.04cm}
    \slidelead{Estimand: selisih antarmode pada \textit{apparatus} ini, bukan efek universal paradigma.}
\end{frame}
}

% =============================================
% 6. Pipeline
% =============================================
{
\usebackgroundtemplate{\backgroundcontent}
\begin{frame}
    \frametitle{Pipeline Pengumpulan dan Analisis}
    \centering
    \includegraphics[width=0.98\textwidth,height=0.32\textheight,keepaspectratio]{pipeline}
    \vspace{0.06cm}
    \begin{columns}[T]
        \begin{column}{0.32\textwidth}
            \begin{beamercolorbox}[wd=\textwidth, sep=0.10cm, center, rounded=true]{structure}
                \scriptsize\bfseries 210 sel
            \end{beamercolorbox}
            \vspace{0.03cm}
            \scriptsize 7 skenario $\times$ 2 mode $\times$ 15 pengulangan: 10 pengulangan blok waktu dan 5 pengulangan blok validitas.
        \end{column}
        \begin{column}{0.32\textwidth}
            \begin{beamercolorbox}[wd=\textwidth, sep=0.10cm, center, rounded=true]{structure}
                \scriptsize\bfseries 207 baris dianalisis
            \end{beamercolorbox}
            \vspace{0.03cm}
            \scriptsize Empat baris mentah dikecualikan; tiga sel tanpa baris analisis; satu sel dianalisis lewat percobaan ulang.
        \end{column}
        \begin{column}{0.32\textwidth}
            \begin{beamercolorbox}[wd=\textwidth, sep=0.10cm, center, rounded=true]{structure}
                \scriptsize\bfseries Analisis
            \end{beamercolorbox}
            \vspace{0.03cm}
            \scriptsize Wilcoxon atau Mann-Whitney, koreksi Holm, Cliff's $\delta$, Kaplan-Meier dan RMST.
        \end{column}
    \end{columns}
    \vspace{0.06cm}
    \smallnote{Bekukan \textit{apparatus}, pemeriksaan kesiapan, penangguhan 60 detik, dan horizon pengamatan 300 detik. Tidak ada imputasi nilai hilang.}
\end{frame}
}

% =============================================
% 7. Drift scenarios
% =============================================
{
\usebackgroundtemplate{\backgroundcontent}
\begin{frame}
    \frametitle{Tujuh Skenario Drift}
    \centering
    \scriptsize
    \begin{tabular}{@{}>{\raggedright\arraybackslash}p{1.25cm}>{\raggedright\arraybackslash}p{5.0cm}>{\raggedright\arraybackslash}p{5.25cm}@{}}
        \toprule
        \textbf{Skenario} & \textbf{Intervensi} & \textbf{Batas interpretasi} \\
        \midrule
        A & Menurunkan replika web ke 1 & Kontrol HPA; replika dikecualikan dari R1 \\
        B & Mengubah tag \textit{image} kontainer web & Drift terkelola pada \textit{field} \textit{image} \\
        C & Mengubah nilai \texttt{INSTANCE\_TAGLINE} pada ConfigMap & Drift terkelola; deklarasi tidak diubah \\
        D & Menghapus Deployment web & Pemulihan dinilai terhadap inventaris dan nilai dasar \\
        E & Mengubah limit memori dari \texttt{2Gi} menjadi \texttt{1536Mi} & Hanya \textit{field} yang ditetapkan yang dinilai \\
        F & Mengubah \texttt{targetPort} Service dari \texttt{http} menjadi \texttt{65535} & Drift Service; dilaporkan terpisah dari keluarga B-E \\
        G & Menambahkan anotasi metadata pada SealedSecret & Kasus deskriptif; mekanisme tidak diatribusikan \\
        \bottomrule
    \end{tabular}
    \vspace{0.10cm}
    \smallnote{A adalah kontrol HPA. F tidak masuk keluarga efikasi B-E. G tidak dipakai sebagai uji \textit{self-heal} ArgoCD.}
\end{frame}
}

% =============================================
% 8. R1 validity
% =============================================
{
\usebackgroundtemplate{\backgroundcontent}
\begin{frame}
    \frametitle{R1: Validitas Pemulihan, Bukan Perbandingan}
    \begin{columns}[T]
        \begin{column}{0.50\textwidth}
            \scriptsize
            \textbf{Empat sumber daya yang diperiksa}
            \begin{itemize}\setlength{\itemsep}{1pt}
                \item \texttt{Deployment/invenio-web}
                \item \texttt{ConfigMap/}\\
                      \texttt{invenio-app-config-6dg2chdbh4}
                \item \texttt{Service/invenio-web}
                \item \texttt{SealedSecret/invenio-app-secrets}
            \end{itemize}
            \vspace{0.06cm}
            \textbf{Tiga titik pengamatan:} kondisi dasar, sampel lengkap pertama setelah injeksi, dan sampel lengkap terakhir.
            \vspace{0.06cm}

            \textbf{\textit{Field} dinamis} seperti \texttt{/spec/replicas} dikecualikan karena dimiliki HPA.
            Kesehatan layanan adalah \textit{guardrail} kesiapan dan \textit{postcondition}; tidak dinilai dan tidak masuk skor.
        \end{column}
        \begin{column}{0.46\textwidth}
            \centering
            \begin{beamercolorbox}[wd=\textwidth, sep=0.24cm, center, rounded=true]{structure}
                \normalsize
                \[
                    R1 = 100
                    \;\Longleftrightarrow\;
                    \text{semua \textit{field} cocok}
                \]
            \end{beamercolorbox}
            \vspace{0.08cm}
            \scriptsize
            Pemeriksa dan evaluator memakai kontrak beku yang sama, sehingga B--F harus kembali ke 100 bila prosedur bekerja.
            Penyebutnya empat sumber daya; satu ketidakcocokan menghasilkan 75.
        \end{column}
    \end{columns}
    \vspace{0.06cm}
    \slidelead{R1 memeriksa apakah manipulasi bekerja; R1 tidak diuji antarmode dan tidak menjadi hasil komparatif.}
\end{frame}
}

% =============================================
% 9. RQ timing
% =============================================
{
\usebackgroundtemplate{\backgroundcontent}
\begin{frame}
    \frametitle{Mengukur Waktu: Tiga Titik Akhir}
    \centering
    \scriptsize
    \begin{tabular}{@{}>{\raggedright\arraybackslash}p{2.1cm}>{\raggedright\arraybackslash}p{9.8cm}@{}}
        \toprule
        \textbf{Kejadian} & \textbf{Definisi operasional} \\
        \midrule
        $t_0$ & Tepat sebelum perintah injeksi pertama dijalankan. \\
        $t_d$ & Selesainya pembacaan atau evaluasi sinyal deteksi pertama yang memenuhi jendela. \\
        $t_e$ & Selesainya evaluasi perbedaan \textit{field} pertama yang mengidentifikasi sumber daya dan \textit{field} berbeda. \\
        $t_r$ & Observasi pertama ketika inventaris dan seluruh \textit{field} R1 kembali cocok, dikonfirmasi observasi berikutnya. \\
        \bottomrule
    \end{tabular}
    \par\vspace{0.16cm}
    \begin{columns}[T]
        \begin{column}{0.48\textwidth}
            \raggedright\scriptsize
            \textbf{Bukti asimetris:} pada mode imperatif $t_e=t_d$ menurut konstruksi; pada GitOps $t_e$ berasal dari kolektor independen.
        \end{column}
        \begin{column}{0.48\textwidth}
            \raggedright\scriptsize
            \textbf{Cadence:} \textit{sampler} 1 detik, pemeriksa imperatif 5 detik; interval rekonsiliasi ArgoCD tidak diukur terpisah.
        \end{column}
    \end{columns}
    \vspace{0.06cm}
    \smallnote{Titik akhir yang tidak tercapai dalam 300 detik dicatat tersensor kanan; waktu hilang tidak diimputasi nol.}
\end{frame}
}

% =============================================
% 10. Data and analysis
% =============================================
{
\usebackgroundtemplate{\backgroundcontent}
\begin{frame}
    \frametitle{Pengumpulan dan Analisis}
    \begin{columns}[T]
        \begin{column}{0.46\textwidth}
            \centering
            \scriptsize
            \begin{tabular}{@{}lr@{}}
                \toprule
                \textbf{Kategori} & \textbf{Jumlah} \\
                \midrule
                Sel direncanakan & 210 \\
                Sel terminal & 210 \\
                Baris mentah & 211 \\
                Baris analisis & 207 \\
                Baris dieksklusi & 4 \\
                \bottomrule
            \end{tabular}
            \par\vspace{0.10cm}
            \raggedright
            193 sel \textit{observed}, 15 tersensor kanan, dan 2 kegagalan instrumentasi.
            Tiga sel tidak menghasilkan baris analisis; satu sel dianalisis lewat percobaan ulang.
            Tidak ada nilai yang diimputasi.
        \end{column}
        \begin{column}{0.50\textwidth}
            \scriptsize
            \textbf{Uji komparatif}
            \begin{itemize}\setlength{\itemsep}{2pt}
                \item Wilcoxon \textit{signed-rank} bila pasangan lengkap; Mann-Whitney U sebagai cadangan eksploratif.
                \item 12 hipotesis B-E untuk deteksi, bukti, dan pemulihan; $\alpha=0{,}05$ dengan koreksi Holm.
                \item Cliff's $\delta$ dengan interval kepercayaan \textit{bootstrap}; Kaplan-Meier dan RMST untuk data tersensor.
            \end{itemize}
            \vspace{0.04cm}
            A, F, dan G dilaporkan di luar keluarga efikasi B-E.
            Kegagalan dan penyensoran dilaporkan terbuka.
        \end{column}
    \end{columns}
    \vspace{0.06cm}
    \slidelead{Tidak ada nilai yang diimputasi; kegagalan dan penyensoran dilaporkan terbuka.}
\end{frame}
}

% =============================================
% 11. Validity result
% =============================================
{
\usebackgroundtemplate{\backgroundcontent}
\begin{frame}
    \frametitle{Hasil Validitas: B--F Konvergen, G Batas}
    \centering
    \scriptsize
    \begin{tabular}{@{}cccc@{}}
        \toprule
        \textbf{Skenario} & \textbf{GitOps} & \textbf{Imperatif} & \textbf{Catatan} \\
        \midrule
        A & 100 / 100 / 100 & 100 / 100 / 100 & Replika dimiliki HPA \\
        B--F & 100 / 75 / 100 & 100 / 75 / 100 & Kembali ke predikat 100 \\
        G & 100 / 75 / \textbf{75} & 100 / 75 / 100 & Anotasi SealedSecret \\
        \bottomrule
    \end{tabular}
    \vspace{0.06cm}

    \raggedleft\scriptsize Format: kondisi dasar / pascainjeksi / pascapemulihan. $n=5$ per skenario per mode.
    \vspace{0.10cm}

    \begin{columns}[T]
        \begin{column}{0.48\textwidth}
            \raggedright\scriptsize
            B--F kembali ke 100 pada kedua mode. Pemeriksa dan evaluator berbagi kontrak beku, sehingga konvergensi adalah syarat validitas manipulasi.
        \end{column}
        \begin{column}{0.48\textwidth}
            \raggedright\scriptsize
            G satu-satunya pengecualian: GitOps tidak mengembalikan anotasi dalam 300 detik. Penyebabnya tidak diatribusikan; G deskriptif, bukan uji \textit{self-heal}.
        \end{column}
    \end{columns}
    \vspace{0.04cm}
    \slidelead{R1 tidak dibandingkan antarmode dan tidak diuji; tabel audit lengkap ada di lampiran.}
\end{frame}
}

% =============================================
% 12. RQ result
% =============================================
{
\usebackgroundtemplate{\backgroundcontent}
\begin{frame}
    \frametitle{Hasil Waktu RQ: B--E}
    \centering
    \scriptsize
    \begin{tabular}{@{}lcccc@{}}
        \toprule
        \textbf{Skenario} & \textbf{Deteksi $\delta$} & \textbf{Bukti $\delta$} & \textbf{Pemulihan $\delta$} & \textbf{$p$ Holm} \\
        \midrule
        B & 1,00 & 1,00 & 0,48 & 0,0176--0,0195 \\
        C & 1,00 & 1,00 & 0,75 & 0,0054--0,0176 \\
        D & 1,00 & 1,00 & 1,00 & 0,0054--0,0176 \\
        E & 1,00 & 1,00 & 0,86 & 0,0176 \\
        \bottomrule
    \end{tabular}
    \vspace{0.06cm}

    \raggedleft\scriptsize
    $n=10$ pasangan kecuali C (8 versus 10), D deteksi (4 versus 10), dan D bukti (8 versus 10).
    \vspace{0.08cm}

    \begin{columns}[T]
        \begin{column}{0.48\textwidth}
            \raggedright\scriptsize
            12 perbandingan B--E: $p$ terkoreksi 0,0054--0,0195 di bawah model uji eksploratif, bukan estimasi populasi. Arah positif berarti imperatif lebih lama.
        \end{column}
        \begin{column}{0.48\textwidth}
            \raggedright\scriptsize
            Pemulihan B: $\delta=0{,}48$ dengan IK 95\% $[-0{,}08;\,1{,}00]$ yang mencakup nol, sehingga bukan keunggulan yang jelas.
        \end{column}
    \end{columns}
    \vspace{0.04cm}
    \slidelead{Deteksi dan bukti paling terpisah; pemulihan tidak seragam.}
\end{frame}
}

% =============================================
% 13. RQ pattern
% =============================================
{
\usebackgroundtemplate{\backgroundcontent}
\begin{frame}
    \frametitle{Cadence dan Bukti yang Asimetris}
    \begin{columns}[T]
        \begin{column}{0.55\textwidth}
            \centering
            \includegraphics[width=\columnwidth,height=0.62\textheight,keepaspectratio]{r2-gaps}
        \end{column}
        \begin{column}{0.41\textwidth}
            \scriptsize
            \textbf{Pola utama}
            \begin{itemize}\setlength{\itemsep}{3pt}
                \item Selisih median deteksi dan bukti B serta E: 4,0 sampai 4,3 detik.
                \item Pemeriksa imperatif 5 detik, \textit{sampler} 1 detik; selisih berdekatan dengan cadence, bukan estimasi murni respons pengendali.
                \item Pemulihan B 2,95 detik dan E 3,03 detik; D tetap 5,05 detik.
                \item Pada mode imperatif $t_e=t_d$ menurut konstruksi.
            \end{itemize}
        \end{column}
    \end{columns}
    \vspace{0.02cm}
    \smallnote{Pola konsisten dengan mekanisme pengamatan, bukan bukti bahwa seluruh selisih berasal dari mekanisme itu.}
\end{frame}
}

% =============================================
% 14. Limits
% =============================================
{
\usebackgroundtemplate{\backgroundcontent}
\begin{frame}
    \frametitle{Keterbatasan yang Membatasi Klaim}
    \centering
    \scriptsize
    \begin{tabular}{@{}>{\raggedright\arraybackslash}p{5.2cm}>{\raggedright\arraybackslash}p{7.2cm}@{}}
        \toprule
        \textbf{Keterbatasan} & \textbf{Dampak pada klaim} \\
        \midrule
        Satu klaster, satu \textit{stack}, satu implementasi GitOps & Tidak dapat digeneralisasi ke semua organisasi atau alat GitOps \\
        Seluruh run GitOps mendahului seluruh run imperatif & Efek periode dan efek urutan tidak dapat dipisahkan dari efek mode \\
        Pengulangan berbagi satu \textit{stack} dan satu periode & Pengulangan pseudoreplikat; $p$ ringkasan deskriptif, bukan estimasi kausal \\
        \textit{Sampler} 1 detik, pemeriksa 5 detik, horizon 300 detik & Ketelitian waktu terbatas pada resolusi instrumen \\
        15 pengulangan per sel & Bukan jaminan daya statistik per skenario \\
        R1 mengukur kesetaraan \textit{field} & Tidak menilai kebenaran semantik atau keamanan konfigurasi \\
        G dan HPA & G tidak dapat digeneralisasi; HPA dibatasi 2--4 replika \\
        \bottomrule
    \end{tabular}
    \vspace{0.08cm}
    \slidelead{Hasil berlaku untuk \textit{apparatus}, protokol, dan cakupan \textit{field} yang benar-benar diuji.}
\end{frame}
}

% =============================================
% 15. Claim boundaries
% =============================================
{
\usebackgroundtemplate{\backgroundcontent}
\begin{frame}
    \frametitle{Batas Klaim: Cara Membaca Hasil}
    \scriptsize
    \begin{itemize}\setlength{\itemsep}{5pt}
        \item \textbf{Urutan blok.} Seluruh run GitOps mendahului seluruh run imperatif; urutan tetap dan tidak diacak ulang. Pengulangan bersifat pseudoreplikat, sehingga nilai $p$ bukan estimasi kausal.
        \item \textbf{Bukti asimetris.} $t_e$ imperatif sama dengan $t_d$ menurut konstruksi; $t_e$ GitOps berasal dari kolektor independen. Perbandingan $t_e$ tidak setara secara mekanisme.
        \item \textbf{Cadence.} Selisih 4,0--4,3 detik berdekatan dengan interval pemeriksa 5 detik dan \textit{sampler} 1 detik, sehingga nilainya bukan estimasi murni respons pengendali.
        \item \textbf{Pemulihan B.} $\delta=0{,}48$ dengan IK 95\% $[-0{,}08;\,1{,}00]$ yang mencakup nol; tidak dibaca sebagai keunggulan yang jelas.
        \item \textbf{Batas G.} G deskriptif dan tidak diatribusikan; tidak ada klaim umum bahwa ArgoCD gagal memulihkan drift.
    \end{itemize}
    \vspace{0.12cm}
    \slidelead{Klaim dibatasi pada sinyal yang teramati di bawah cadence masing-masing.}
\end{frame}
}

% =============================================
% 16. Conclusion
% =============================================
{
\usebackgroundtemplate{\backgroundcontent}
\begin{frame}
    \frametitle{Kesimpulan: Waktu RQ dan Validitas R1}
    \scriptsize
    \begin{itemize}\setlength{\itemsep}{4pt}
        \item \textbf{Validitas.} Median R1 untuk B--F kembali ke 100 pada kedua mode; G GitOps tetap 75. B--F adalah syarat validitas manipulasi, bukan hasil komparatif.
        \item \textbf{Waktu.} Pada B--E, GitOps lebih cepat pada sinyal yang teramati di bawah cadence masing-masing untuk 12 perbandingan, dengan $p$ terkoreksi 0,0054--0,0195 di bawah model uji eksploratif.
        \item \textbf{Pemulihan.} Selisih pemulihan tidak seragam: B $\delta=0{,}48$ dengan IK yang mencakup nol, C $0{,}75$, D $1{,}00$, dan E $0{,}86$.
        \item \textbf{Batas.} Tidak ada klaim superioritas umum, kelayakan adopsi, atau kesetaraan semantik konfigurasi. Kecepatan deteksi bukan pengganti jalur pemulihan yang teruji.
    \end{itemize}
    \vspace{0.12cm}
    \slidelead{Satu pertanyaan waktu, satu pemeriksaan validitas, dan batas klaim yang eksplisit.}
\end{frame}
}

% =============================================
% 17. Future work
% =============================================
{
\usebackgroundtemplate{\backgroundcontent}
\begin{frame}
    \frametitle{Pekerjaan Lanjutan}
    \begin{columns}[T]
        \begin{column}{0.48\textwidth}
            \scriptsize
            \textbf{Memisahkan variabel}
            \begin{itemize}\setlength{\itemsep}{4pt}
                \item Replikasi pada tiga \textit{stack} terisolasi untuk memisahkan efek urutan dari efek mode.
                \item Sampling sub-detik dan jam presisi untuk memisahkan mekanisme pengamatan dari ketanggapan pengendali.
                \item Menambah pengulangan berdasarkan daya yang teramati.
            \end{itemize}
        \end{column}
        \begin{column}{0.48\textwidth}
            \scriptsize
            \textbf{Memperluas cakupan}
            \begin{itemize}\setlength{\itemsep}{4pt}
                \item Beban pengguna yang realistis dan pengukuran dampak layanan.
                \item Perbandingan beberapa implementasi GitOps, misalnya Flux dan ArgoCD.
                \item Infrastruktur cadangan independen dan skenario drift berganda.
            \end{itemize}
        \end{column}
    \end{columns}
    \vspace{0.16cm}
    \slidelead{Prioritas: memisahkan variabel yang belum dapat dipisahkan, bukan menambah pengamatan pada \textit{apparatus} yang sama.}
\end{frame}
}

% =============================================
% 18. Closing
% =============================================
{
\usebackgroundtemplate{\backgroundcontent}
\begin{frame}
    \frametitle{Terima Kasih}
    \vspace{0.16cm}
    \begin{enumerate}\setlength{\itemsep}{7pt}
        \item \textbf{Perbandingan terkontrol pada \textit{workload} multi-service nyata:} satu \textit{stack} InvenioRDM, dua mode berurutan, tujuh skenario drift.
        \item \textbf{Protokol yang dapat diaudit:} definisi waktu dan validitas R1; penanganan data tersensor; keluarga pengujian yang ditetapkan sebelum pengumpulan.
        \item \textbf{Hasil yang dibatasi:} B--F kembali ke predikat R1, dengan G sebagai batas deskriptif; deteksi cepat tidak selalu diikuti pemulihan cepat.
    \end{enumerate}
    \vspace{0.18cm}
    \begin{beamercolorbox}[wd=0.92\textwidth, sep=0.20cm, center, rounded=true]{structure}
        \normalsize\bfseries Diskusi dan pertanyaan
    \end{beamercolorbox}
\end{frame}
}

% =============================================
% BACKUP SLIDES
% =============================================
\appendix

% ---------------------------------------------
% B1: Why only timing and validity
% ---------------------------------------------
{
\usebackgroundtemplate{\backgroundcontent}
\begin{frame}
    \frametitle{Backup: Batas Pengukuran}
    \small
    \begin{itemize}\setlength{\itemsep}{6pt}
        \item Atribusi pelaku mutasi memerlukan audit API server untuk menghubungkan perubahan dengan pelakunya.
        \item Audit API server tidak tersedia pada lingkungan penelitian, sehingga atribusi tidak dapat diukur.
        \item Memaksakan metrik atribusi akan menghasilkan klaim yang tidak dapat diverifikasi. Karena itu, atribusi dilaporkan sebagai keterbatasan.
        \item Waktu dan validitas R1 tetap dapat diukur dengan aturan yang sama pada kedua mode dan didefinisikan sebelum pengumpulan.
    \end{itemize}
    \vspace{0.12cm}
    \slidelead{Ketiadaan atribusi adalah batas \textit{apparatus}, bukan hasil yang disembunyikan.}
\end{frame}
}

% ---------------------------------------------
% B2: Cell accounting
% ---------------------------------------------
{
\usebackgroundtemplate{\backgroundcontent}
\begin{frame}
    \frametitle{Backup: Akuntansi Sel dan Data Tersensor}
    \centering
    \scriptsize
    \begin{tabular}{@{}lr>{\raggedright\arraybackslash}p{8.2cm}@{}}
        \toprule
        \textbf{Kategori} & \textbf{Jumlah} & \textbf{Komposisi} \\
        \midrule
        Sel direncanakan & 210 & 7 skenario $\times$ 2 mode $\times$ 15 pengulangan \\
        Sel terminal & 210 & 193 \textit{observed}, 15 tersensor kanan, 2 kegagalan instrumentasi \\
        Baris mentah & 211 & 210 sel ditambah satu percobaan \textit{baseline} yang gagal \\
        Baris analisis & 207 & 192 \textit{observed}, 15 tersensor kanan \\
        Baris dieksklusi & 4 & 2 kegagalan instrumentasi, 1 \textit{batch stop}, 1 \textit{baseline failure} \\
        \bottomrule
    \end{tabular}
    \par\vspace{0.30cm}
    \begin{columns}[T]
        \begin{column}{0.48\textwidth}
            \raggedright\scriptsize
            Tiga sel tidak menghasilkan baris analisis. Satu sel, \texttt{gitops-C-r08}, dianalisis melalui percobaan ulang setelah kegagalan \textit{baseline} pra-injeksi.
        \end{column}
        \begin{column}{0.48\textwidth}
            \raggedright\scriptsize
            Seluruh 15 kejadian sensoring pemulihan berasal dari mode GitOps pada skenario G. Tidak ada nilai tersensor yang diubah menjadi kegagalan atau nol.
        \end{column}
    \end{columns}
\end{frame}
}

% ---------------------------------------------
% B3: Scenario G
% ---------------------------------------------
{
\usebackgroundtemplate{\backgroundcontent}
\begin{frame}
    \frametitle{Backup: Mengapa Skenario G Deskriptif?}
    \small
    \begin{itemize}\setlength{\itemsep}{6pt}
        \item Intervensi G adalah anotasi metadata pada SealedSecret; nilai \texttt{spec.encryptedData} tidak dibaca atau ditulis.
        \item Mode imperatif memulihkan anotasi melalui pemeriksa generik. Mode GitOps tidak memulihkannya dalam 300 detik.
        \item Keluaran baris mencatat \texttt{recovery\_controller} sebagai \texttt{unattributed}.
        \item Penelitian ini tidak menyimpulkan apakah penyebabnya cakupan \textit{field} ArgoCD, perilaku pengendali SealedSecret, atau faktor lain.
        \item Karena mekanismenya tidak teratribusi, G tidak dipakai sebagai uji \textit{self-heal} dan tidak digabungkan ke keluarga B-E.
    \end{itemize}
\end{frame}
}

% ---------------------------------------------
% B4: Main references
% ---------------------------------------------
{
\usebackgroundtemplate{\backgroundcontent}
\begin{frame}
    \frametitle{Backup: Rujukan Utama}
    \scriptsize
    \begin{itemize}\setlength{\itemsep}{4pt}
        \item Kagganti Nataraja, P. (2025). \textit{A Security-Centric Analysis of Declarative and Imperative Deployment Approaches in Kubernetes-Based Application Environments}. National College of Ireland.
        \item Matubber, M. E. (2025). \textit{Managing 5G Kubernetes Infrastructures Using GitOps Principles}. KTH Royal Institute of Technology.
        \item Pohjola, O. (2025). \textit{Mitigating Configuration Drift in Infrastructure-as-Code Systems}. Aalto University.
        \item Rahman, A., Shamim, S. I., Bose, D. B., dkk. (2023). Security Misconfigurations in Open Source Kubernetes Manifests: An Empirical Study. \textit{ACM TOSEM}.
        \item Shrestha, R. dan Ali, A. A. N. (2024). Configuration Management in Kubernetes Environments: A GitOps Approach. \textit{IEEE/ACM UCC}.
        \item Zhang, Y., Paul, U., d'Amorim, M., dan Rahman, A. (2026). Configuration Defects in Kubernetes. \textit{IEEE TSE}, 52(3), 1125-1144.
    \end{itemize}
\end{frame}
}

\end{document}
